\section{Comparison of Two Treatment Group with Continuous Endpoints}



\subsection{Hypothesis}
각 집단의 크기를 $n_i$이고, 집단의 갯수가 $k$개라면 검정통계량 $F$는 다음과 같이 정의한다. 
    \begin{equation*}
        F = \frac{MS_{between}}{MS_{within}} = \frac{SSB/(k-1)}{SSW/(N-k)}
    \end{equation*}
가설검정은 다음과 같이 진행된다. 
\begin{itemize}[leftmargin=2.5em]
    \item 귀무가설이 참일때 F통계량은 자유도 $(k-1, N-k)$인 중심-F분포를 따른다. 
    \item 대립가설이 참일때 F통계량은 자유도 $(k-1, N-k, \lambda)$인 비중심-F분포를 따른다. 
\end{itemize}
이 이론에 따라, 비중심 모수라는 것이 무엇인지에 대한 검토가 필요하다. 






\subsection{Non-centrality Parameter}
검정력 계산을 위해서는 대립가설 조건에서 F-분포가 오른쪽으로 얼마나 이동하는지를 나타내는 비중심 모수가 필요하다. 
\\
    \begin{equation*}
        \lambda = \frac{\sum^k_{i=1}n_i(\mu_i-\mu)^2}{\sigma^2}
    \end{equation*}
\\
\noindent
모든 집단의 표본의 크기 $n$으로 동일 균형 디자인의 경우 $\lambda$는 \uline{\tb{효과 크기 지표}}인 cohen's $f$를 사용하여 표현할 수 있다. 여기서 $\sigma_\mu$는 각 그룹 평균들의 표준편차이다. $\sigma$는 그룹내 공통 표준편차 (잔차 표준편차)이다. 
\\
    \begin{equation*}
        \lambda = N\cdot f^2, \qquad f = \frac{\sigma_\mu}{\sigma} = \sqrt{\frac{\frac{1}{k}\sum{(\mu_{i}-\mu)^2}}{\sigma^2}}
    \end{equation*}




    
\subsection{Power Calculation}
\label{power_calculation}
검정력은 귀무가설 하에서 설정한 유의수준 $\alpha$의 기각치를 사용하여 비중심 $F$분포의 확률밀도함수를 적분해서 얻는다. 
    \begin{align*}
        1-\beta &= P(F(k-1, N-k, \lambda) \geq F_{crit}\\ &= \int^\infty_{F_{crit}}P_{nc}(x; k-1, N-k, \lambda)dx
    \end{align*}
비중심모수를 알기 위해서는 $\sigma$에 대한 정보가 필요하다. 그런데 실무에서는 모집단에 관한 정보를 정확히 얻기 힘들다. 그래서 효과크기 $f$에 일반적인 기준이 존재한다. 
\begin{tcolorbox}[arc=4 mm, colback=blue!5, colframe=black, boxrule=0.1 mm, breakable]
    \begin{itemize}[leftmargin=1.5em]
        \item $f=0.10$: 작은효과에 해당한다. 집단간 차이가 크지 않아서 육안으로 식별하기 어렵다. 유의미한 검정을 위해서는 대표본이 요구된다. 
        \item $f=0.25$: \uline{일상적인 통계 연구에서} 가장 흔하게 가정하는 수준이다. 
        \item $f=0.40$: 집단간 차이가 명확하고, 작은 표본으로 높은 검정력을 확보할 수 있다. 
    \end{itemize}
\end{tcolorbox}






\subsection{Sample Size Calculation}
정확한 해를 구할 수 있는 방정식은 없다. 비중심 F분포의 누적분포함수의 역함수를 이용해서 반복 계산하는 방식을 취한다. 해당 알고리즘을 정리해보자. 
\begin{tcolorbox}[arc=4 mm, colback=green!5, colframe=black, boxrule=0.1 mm, breakable]
    \begin{itemize}[leftmargin = 5.5 em]
        \item[STEP. 01] \tb{효과의 크기 계산} : Cohen's $f$ 값을 정한다. 연구 데이터의 예측치나 선행 연구를 기반으로 한다. 연구자료가 없을 경우는 일반적 기준을 적용한다. 세부 기준은 \ref{power_calculation}의 내용을 참고한다. 
        \item[STEP. 02] \tb{기각치 $(F_{crit})$ 결정} : 주어진 유의수준 $(\alpha)$와 그룹수 $(k)$, 총샘플사이즈 $(N)$을 가정하고, 중심 F분포를 이용해서 기각치를 계산한다. 
        \item[STEP. 03] \tb{비중심모수 계산} : 가정한 샘플사이즈$(N)$와 효과크기 $(f)$를 사용해서 비중심모수 $(\lambda)$를 계산한다. 
            \begin{equation*}
                \lambda = N\cdot f^2
            \end{equation*}
        \item[STEP. 04] \tb{실제 검정력 계산 및 비교}: 비중심 F분포에서 기각치보다 큰 영역의 확률 (검정력)계산한다. 
            \begin{equation*}
                \text{Calculated Power} = P(F_{nc} > F_{crit})
            \end{equation*}
        \item[STEP. 05] \tb{반복 수정} : 여기서는 다음 기준으로 알고리즘 진행여부를 결정한다. 
            \begin{itemize}[nosep]
                \item [(1)] $\text{Calculated Power} < \text{Targeted Power}$: $N$을 증가시킨 후 STEP.02로 이동
                \item [(2)] $\text{Calculated Power} \geq \text{Targeted Power}$: 현재 $N$을 최종 sample size로 결정
            \end{itemize}
        \item[STEP. 06] \tb{그룹별 샘플 사이즈 할당} : 각 그룹의 샘플사이즈가 동일하다고 가정한다. 그러면 각 그룹당 샘플 사이즈는 다음과 같다. 
            \begin{equation*}
                n_i = \frac{N}{k}
            \end{equation*}
    \end{itemize}
\end{tcolorbox}
\noindent
상기 알고리즘에서 고민해 봐야할 부분은 $N$을 증가시키는 단계 별 수준이다. 여기에 대해서 고찰해 보자. 어느 정도 만큼 $N$을 증가시키면 될까? 현재까지 적용하고 있는 방식은 크게 3가지가 있다. \tb{균형설계를 유지하는 논리방식이 가장 적절}하다. 
\begin{itemize}[leftmargin=2.5em]
    \item[(1)] \uline{\textbf{균형설계의 유지}}: 현재 설계 논리상 각 집단은 동일한 표본 수로 구성된다. 즉 모든 집단이 동일한 샘플 사이즈를 가진 균형설계를 가정하고 있다. 수리적 구조상 $N$은 $n_i$의 $k$배 형태로 증가한다. 총 샘플사이즈가 $k$개 증가하기 때문에 \tb{각 집단에 1명씩 추가로 할당하는 형태}가 된다. 
        \begin{equation*}
            N^* = N + k
        \end{equation*}
    \item[(2)] \uline{\textbf{검정력의 단조 증가성}} : 검정력은 표본의 크기 $N$이 커질수록 항상 증가하는 단조증가 형태를 띤다. 따라서, 샘플사이즈를 조금씩 증가시키면서 목표 검정력을 만족하는 최소 샘플사이즈를 도출하는 것이 좋다. 
    \item[(3)] \uline{\textbf{소프트웨어 내부의 탐색 알고리즘 논리}} : 난점은 상기 논리가 계산 차원에서는 시간이 많이 들 수 있다. 따라서, 하기 방식을 이용하여 목표치를 도출하는 방식이 좋다. 
\end{itemize}

\begin{tcolorbox}[arc=4mm, colback=blue!5, colframe=black, boxrule=0.1 mm]
    \begin{itemize}[leftmargin = 1em, nosep]
        \item 이분법 (Bisection method), 뉴턴-랩슨법 (Newton-Raphson method)
        \item 상기 방식으로 목표치에 도달하면 집단수 $k$만큼 조금씩 샘플사이즈를 증가시키면서 목표검정력에 도달하는 최소 샘플수를 확정한다. 
    \end{itemize}
\end{tcolorbox}
